\chapter{Requerimientos funcionales, no funcionales y trazabilidad}
\label{chap:requerimientos}

\section{Propósito del capítulo}

Este capítulo formaliza qué debe hacer el prototipo, bajo qué condiciones y cómo se verifica cada capacidad, mediante la cadena necesidad, historia de usuario, requerimiento, componente, criterio de aceptación y evidencia. Los identificadores RF, RNF y RV conservan el significado establecido en el Capítulo 8 y se amplían para cubrir autenticación, persistencia, trazabilidad distribuida y compatibilidad volumétrica.
% CROSS_REFERENCE_MAPPING_REQUIRED

\section{Fuentes y criterios de formalización}

La especificación adopta la distinción habitual entre requerimientos funcionales, que describen qué debe hacer el sistema, y no funcionales, que describen bajo qué condiciones de calidad debe hacerlo \parencite{wiegers_software_requirements_2013}. Se construye a partir de cuatro grupos de fuentes:

\begin{itemize}
  \item User research: necesidades N-01 a N-12 y decisiones derivadas de las entrevistas EN-01 a EN-05.
  \item Alcance del producto: separación sagital–axial, revisión humana, salida no diagnóstica, ingesta de estudios DICOM completos (con MHA/NIfTI como compatibilidad técnica) y exclusiones clínicas.
  \item Implementación real: Frontend, Backend y AI Module, junto con PostgreSQL, assets durables, contratos y pruebas.
  \item Roadmap técnico: P9, P10, P10.5 y evolución posterior hacia Google Cloud y comparación longitudinal descriptiva.
\end{itemize}

No se incorporan requerimientos únicamente para completar una cantidad. Cada elemento debe estar vinculado con una necesidad, una capacidad comprobable o una decisión explícita de alcance. Cuando una función todavía no está implementada, se indica como «En curso» o «Futuro»; no se la presenta como terminada.

\section{Actores y responsabilidades}

{\small
\setlength{\LTleft}{0pt}
\setlength{\LTright}{0pt}
\setlength{\tabcolsep}{2pt}
\renewcommand{\arraystretch}{1.18}
\begin{longtable}{>{\raggedright\arraybackslash}p{0.12\textwidth} >{\raggedright\arraybackslash}p{0.28\textwidth} >{\raggedright\arraybackslash}p{0.52\textwidth}}
\caption{Actores y responsabilidades}\label{tab:req-actores}\\
\hline
\textbf{Código} & \textbf{Actor} & \textbf{Responsabilidad principal} \\
\hline
\endfirsthead
\caption[]{Actores y responsabilidades (continuación)}\\
\hline
\textbf{Código} & \textbf{Actor} & \textbf{Responsabilidad principal} \\
\hline
\endhead
\hline
\endfoot
A1 & Profesional revisor — roles DOCTOR o REVIEWER & Crear o consultar estudios, ejecutar análisis, inspeccionar imágenes y resultados, revisar mediciones y registrar una decisión profesional. \\
A2 & Administrador — rol ADMIN & Gestionar usuarios y permisos, consultar diagnósticos técnicos y supervisar la disponibilidad operativa del sistema. \\
A3 & Backend de producto & Validar autenticación y autorización, aplicar reglas de negocio, persistir estudios y corridas, publicar assets autorizados y orquestar la comunicación con el AI Module. \\
A4 & AI Module & Validar modelos y artifacts, preprocesar volúmenes, ejecutar inferencia, generar máscaras, overlays, mediciones, metadata y artifacts técnicos. \\
A5 & Equipo de desarrollo e investigación & Mantener datasets, modelos, manifests, contratos, pruebas, evidencia y documentación reproducible. \\
\hline
\end{longtable}
}

El profesional es el actor humano principal. El Backend y el AI Module se documentan como actores de sistema para hacer explícitas las responsabilidades distribuidas y evitar atribuir al Frontend capacidades que dependen de otras capas.

\section{Estados de requerimiento y criterio de prioridad}

{\small
\setlength{\LTleft}{0pt}
\setlength{\LTright}{0pt}
\setlength{\tabcolsep}{2pt}
\renewcommand{\arraystretch}{1.18}
\begin{longtable}{>{\raggedright\arraybackslash}p{0.22\textwidth} >{\raggedright\arraybackslash}p{0.70\textwidth}}
\caption{Estados de requerimiento y criterio de prioridad}\label{tab:req-estados}\\
\hline
\textbf{Estado} & \textbf{Definición} \\
\hline
\endfirsthead
\caption[]{Estados de requerimiento y criterio de prioridad (continuación)}\\
\hline
\textbf{Estado} & \textbf{Definición} \\
\hline
\endhead
\hline
\endfoot
Implementado & La capacidad existe en el producto y dispone de evidencia técnica o funcional. \\
Parcial & Existe una parte funcional, pero falta completar el alcance definido o su validación end-to-end. \\
En curso & El contrato o desarrollo está iniciado y posee tareas bloqueantes identificadas. \\
Futuro & La capacidad se conserva como evolución y no integra el alcance inmediato. \\
\hline
\end{longtable}
}

La prioridad se clasifica como alta, media o baja. Una prioridad alta corresponde a seguridad, flujo principal, revisión humana, trazabilidad o funciones que bloquean otras capas. Una prioridad media identifica capacidades relevantes que no bloquean el recorrido principal. Una prioridad baja corresponde a evoluciones posteriores, como comparación longitudinal avanzada.

\section{Flujos alternativos y excepciones}

FA-01 — AI Module no disponible. El Backend conserva su disponibilidad, devuelve un estado degradado explícito y evita presentar una respuesta contractual como inferencia real. Los estudios persistidos continúan siendo consultables dentro del alcance de los assets ya almacenados.

FA-02 — Archivo inválido o incompatible. La carga se rechaza con un mensaje comprensible y no se crea una corrida parcial presentada como válida.

FA-03 — Plano axial ausente. El sistema debe informar la ausencia y permitir un flujo sagital cuando el contrato lo admita, sin fabricar resultados axiales.

FA-04 — Modelo o artifact no disponible. La corrida no se marca como inferencia real. El diagnóstico técnico debe indicar qué modelo o artifact falta.

FA-05 — Usuario sin permisos. El Backend rechaza la operación y el Frontend no expone acciones para las que el rol no está autorizado.

FA-06 — Corrida histórica. Si el estudio solo conserva input.png y overlay.png, debe poder abrirse en modo compatible, sin inventar un catálogo de slices inexistente.

FA-07 — Corte sin resultados. El visor muestra la imagen base y el estado «Sin resultados automáticos en este corte».

FA-08 — Asset faltante o corrupto. El sistema informa la ausencia, conserva la trazabilidad del error y evita reemplazarlo por un contenido engañoso.

\section{Historias de usuario priorizadas}

{\footnotesize
\setlength{\LTleft}{0pt}
\setlength{\LTright}{0pt}
\setlength{\tabcolsep}{2pt}
\renewcommand{\arraystretch}{1.18}
\begin{longtable}{>{\raggedright\arraybackslash}p{0.10\textwidth} >{\raggedright\arraybackslash}p{0.50\textwidth} >{\raggedright\arraybackslash}p{0.14\textwidth} >{\raggedright\arraybackslash}p{0.18\textwidth}}
\caption{Historias de usuario priorizadas}\label{tab:req-historias-usuario}\\
\hline
\textbf{ID} & \textbf{Historia de usuario} & \textbf{Prioridad} & \textbf{Estado} \\
\hline
\endfirsthead
\caption[]{Historias de usuario priorizadas (continuación)}\\
\hline
\textbf{ID} & \textbf{Historia de usuario} & \textbf{Prioridad} & \textbf{Estado} \\
\hline
\endhead
\hline
\endfoot
HU-01 & Como usuario autorizado, quiero iniciar y cerrar sesión para acceder únicamente a las funciones permitidas por mi cuenta. & Alta & Implementado \\
HU-02 & Como profesional, quiero consultar una worklist de estudios para identificar casos disponibles, pendientes y revisados. & Alta & Implementado \\
HU-03 & Como profesional, quiero crear un estudio asociado a una referencia de sujeto de-identificada para organizar casos sin almacenar identidad clínica innecesaria. & Alta & Implementado \\
HU-04 & Como profesional, quiero cargar un estudio DICOM completo y utilizar las series compatibles identificadas por el sistema para ejecutar el análisis. & Alta & Implementado \\
HU-05 & Como profesional, quiero ejecutar el análisis y conocer el estado de la corrida para distinguir procesamiento, resultado, degradación y error. & Alta & Implementado \\
HU-06 & Como profesional, quiero observar la imagen original con el overlay automático para comprobar visualmente la segmentación. & Alta & Implementado \\
HU-07 & Como profesional, quiero recorrer todos los cortes del volumen para analizar el contexto alrededor del corte inferido. & Alta & Implementado; integración final del contrato P10.9 en Frontend pendiente. \\
HU-08 & Como profesional, quiero consultar mediciones vinculadas con estructura, plano y corte para entender de dónde surge cada valor. & Alta & Implementado \\
HU-09 & Como profesional, quiero aceptar, observar, corregir o descartar resultados para mantener la revisión humana obligatoria. & Alta & Implementado \\
HU-10 & Como profesional, quiero reabrir un estudio ya procesado sin repetir la inferencia para conservar continuidad y evitar cómputo innecesario. & Alta & Implementado; recorrido final con interfaz sujeto a integración P10.9 del Frontend. \\
HU-11 & Como usuario, quiero recibir mensajes claros cuando la IA no está disponible o falta información para no confundir una degradación con un resultado válido. & Alta & Parcial \\
HU-12 & Como administrador, quiero consultar salud, modelos y configuración operativa para diagnosticar fallas del sistema. & Media & Implementado \\
HU-13 & Como administrador, quiero gestionar activación y roles de usuarios para controlar el acceso al producto. & Alta & Implementado \\
HU-14 & Como profesional, quiero comparar mediciones de estudios previos del mismo sujeto para observar cambios descriptivos en el tiempo. & Baja & Futuro \\
\hline
\end{longtable}
}

\section{Requerimientos funcionales}

{\scriptsize
\setlength{\LTleft}{0pt}
\setlength{\LTright}{0pt}
\setlength{\tabcolsep}{2pt}
\renewcommand{\arraystretch}{1.18}
\begin{longtable}{>{\raggedright\arraybackslash}p{0.08\textwidth} >{\raggedright\arraybackslash}p{0.31\textwidth} >{\raggedright\arraybackslash}p{0.35\textwidth} >{\raggedright\arraybackslash}p{0.18\textwidth}}
\caption{Requerimientos funcionales}\label{tab:req-funcionales}\\
\hline
\textbf{ID} & \textbf{Requerimiento} & \textbf{Criterio de aceptación} & \textbf{Prioridad / estado / fuente} \\
\hline
\endfirsthead
\caption[]{Requerimientos funcionales (continuación)}\\
\hline
\textbf{ID} & \textbf{Requerimiento} & \textbf{Criterio de aceptación} & \textbf{Prioridad / estado / fuente} \\
\hline
\endhead
\hline
\endfoot
RF-01 & Ingerir un estudio de RM lumbar en formato DICOM (empaquetado en ZIP), identificar sus series, clasificar plano y ponderación, y registrar los inputs compatibles para el análisis. & El sistema recibe un estudio DICOM completo, identifica las series (por ejemplo Sagital T1, Sagital T2/STIR y Axial T2), registra los inputs compatibles y presenta al menos el corte seleccionado por serie. Se mantiene compatibilidad secundaria con inputs individuales MHA/NIfTI para investigación. & Prioridad: Alta Estado: Implementado Fuente: N-03, N-07 \\
RF-02 & Mostrar máscaras o contornos superpuestos sobre la imagen original. & El usuario puede activar y desactivar el overlay y comprobar que corresponde al plano y corte informado. & Prioridad: Alta Estado: Implementado Fuente: N-01; HU-06 \\
RF-03 & Permitir revisión explícita de resultados automáticos. & Cada resultado puede quedar aceptado, observado, corregido o descartado mediante una acción autenticada. & Prioridad: Alta Estado: Implementado Fuente: N-02, N-05; HU-09 \\
RF-04 & Generar mediciones geométricas trazables por estructura, plano, máscara y corte. & Cada medición conserva valor, unidad, estructura o label, plano, sliceIndex y referencia al resultado que la originó. & Prioridad: Alta Estado: Implementado Fuente: N-04; HU-08 \\
RF-05 & Generar una salida estructurada editable y no diagnóstica. & La salida reúne estructuras, mediciones, observaciones, advertencias y revisión, sin cerrar diagnóstico ni recomendar tratamiento. & Prioridad: Alta Estado: Implementado Fuente: N-05, N-07; HU-08, HU-09 \\
RF-06 & Registrar discrepancias, correcciones u observaciones del usuario. & El sistema persiste al menos el estado de revisión, comentario y usuario responsable vinculados con la corrida. & Prioridad: Media Estado: Implementado Fuente: N-08; HU-09 \\
RF-07 & Incorporar un módulo axial complementario sin fusionar artificialmente las taxonomías de los datasets. & La corrida multiplanar ejecuta o registra un resultado axial diferenciado, con modelo, clases y limitaciones propias. & Prioridad: Alta Estado: Implementado Fuente: N-03; HU-04, HU-05 \\
RF-08 & Documentar y habilitar progresivamente la comparación con estudios previos del mismo subjectRef. & La evolución longitudinal se mantiene separada del MVP inmediato y define datos, comparabilidad, deltas descriptivos y revisión profesional. & Prioridad: Baja Estado: Futuro Fuente: N-10; HU-14 \\
RF-09 & Permitir la navegación ordenada de todos los cortes del volumen MHA. & El usuario puede recorrer series y cortes del estudio mediante scroll, teclado y modo cine, con controles de imagen (window/level, zoom) e indicador de corte actual/total; los cortes sin inferencia se muestran como imagen base. & Prioridad: Alta Estado: Implementado (Backend + AI en full-series; visor navegable en Frontend). Integración final del contrato P10.9 con el Frontend, pendiente. Fuente: N-12 \\
RF-10 & Comunicar planos, cortes, modelos, assets o resultados ausentes sin completar artificialmente la información. & El Frontend muestra un estado explícito y el contrato diferencia ausencia, degradación y error. & Prioridad: Alta Estado: Parcial Fuente: N-11; HU-11 \\
RF-11 & Autenticar usuarios y administrar una sesión renovable. & Los endpoints protegidos requieren credenciales válidas, la sesión puede renovarse y el cierre de sesión invalida el acceso desde la interfaz. & Prioridad: Alta Estado: Implementado Fuente: Alcance de seguridad; HU-01 \\
RF-12 & Aplicar autorización por rol y estado de cuenta. & Las acciones de administración, revisión y consulta se habilitan o rechazan según ADMIN, DOCTOR, REVIEWER y estado de aprobación. & Prioridad: Alta Estado: Implementado Fuente: Alcance de seguridad; HU-01, HU-13 \\
RF-13 & Proveer una worklist de estudios y corridas. & El usuario autorizado puede listar estudios, abrir un caso y distinguir estados relevantes del flujo. & Prioridad: Alta Estado: Implementado Fuente: HU-02 \\
RF-14 & Asociar cada estudio con caseId o studyId y una referencia de sujeto de-identificada. & La creación y consulta conservan identificadores estables sin exigir nombre, DNI u otro dato identificatorio. & Prioridad: Alta Estado: Implementado Fuente: N-10; HU-03 \\
RF-15 & Orquestar una corrida multiplanar reconociendo entradas independientes: Sagital T1, Sagital T2 y Axial. & Cada entrada (sagittalT1, sagittalT2, axial) conserva su propio inputId y segmentationRunId; T1 y T2 se procesan de forma independiente y ambas pueden alimentar la línea de hallazgos P10.7. Los resultados se agrupan bajo un identificador común sin asumir correspondencia espacial automática entre planos. & Prioridad: Alta Estado: Implementado Fuente: N-03 \\
RF-16 & Conservar trazabilidad de modelo, artifact y modo de inferencia. & La corrida registra modelKey, versión, hash del artifact, requestedInferenceMode, effectiveInferenceMode y traceId cuando corresponda. & Prioridad: Alta Estado: Implementado Fuente: N-04, N-06, N-08 \\
RF-17 & Persistir resultados, mediciones, revisión y referencias a assets. & El estudio puede recuperarse desde PostgreSQL y storage durable sin depender de la memoria de una sesión del navegador. & Prioridad: Alta Estado: Implementado Fuente: HU-10 \\
RF-18 & Reabrir estudios persistidos cuando el AI Module no está disponible. & El resultado de una corrida se persiste en PostgreSQL y puede reabrirse extremo a extremo, incluyendo mediciones, correcciones profesionales y estado de revisión. & Prioridad: Alta Estado: Implementado (validado en Backend/AI; recorrido con interfaz sujeto a la integración P10.9 del Frontend) Fuente: HU-10 \\
RF-19 & Exponer diagnóstico técnico de salud y modelos. & Un usuario autorizado puede consultar disponibilidad del Backend, AI Module, modelos registrados y estado degradado sin revelar secretos. & Prioridad: Media Estado: Implementado Fuente: HU-12 \\
RF-20 & Servir assets mediante URLs controladas y autorización del Backend. & El navegador no recibe paths internos del AI Module y los assets protegidos requieren una sesión válida. & Prioridad: Alta Estado: Implementado Fuente: RNF-06, RNF-07 \\
RF-21 & Mantener compatibilidad con corridas históricas y contratos versionados. & Un estudio legacy con assets limitados puede abrirse sin inventar slices; los consumidores diferencian versiones y campos opcionales. & Prioridad: Media Estado: En curso Fuente: Contrato P10.5-A; FA-06 \\
RF-22 & Generar hallazgos degenerativos asistivos por nivel (línea P10.7). & El sistema puede generar los ocho tipos de hallazgo definidos por contrato (pfirrmann\_grade, modic\_change, upper\_endplate\_change, lower\_endplate\_change, spondylolisthesis, disc\_herniation, disc\_narrowing, disc\_bulging), asociados al nivel correspondiente, con classification.label y deploymentStatus. No constituye diagnóstico autónomo, requiere revisión profesional y no expone probabilidades al usuario. & Prioridad: Alta Estado: Implementado (Backend + AI) Fuente: estado del arte; P10.7 \\
RF-23 & Gestionar Sagital T1 y Sagital T2 como modalidades independientes. & Cada modalidad conserva inputId y segmentationRunId propios; no se asume correspondencia espacial implícita entre T1 y T2; ambas pueden participar de la línea de hallazgos P10.7. No existe registro automático T1\ensuremath{\leftrightarrow}T2 validado. & Prioridad: Alta Estado: Implementado Fuente: pipeline P10.9 \\
RF-24 & Exponer el resultado de IA de forma segura en el contrato público. & El contrato público no expone probabilidades, logits, salidas crudas, mapas de confianza ni rutas internas del modelo; preserva los campos necesarios para la revisión (classification.label, level, provenance, deploymentStatus y estado de revisión). & Prioridad: Alta Estado: Implementado Fuente: contrato P10.9 \\
RF-25 & Generar un preinforme automático revisable a partir de las estructuras segmentadas y de las mediciones derivadas, editable por el profesional antes de su uso. & El sistema produce un texto estructurado por estructura y nivel, con las mediciones y sus unidades, que el profesional puede modificar, completar o descartar. El preinforme indica de forma explícita que no constituye un informe clínico y conserva el estado de revisión asociado. & Prioridad: Alta Estado: Implementado Fuente: N-04, N-08; HU-09 \\
\hline
\end{longtable}
}

\section{Requerimientos no funcionales}

{\scriptsize
\setlength{\LTleft}{0pt}
\setlength{\LTright}{0pt}
\setlength{\tabcolsep}{2pt}
\renewcommand{\arraystretch}{1.18}
\begin{longtable}{>{\raggedright\arraybackslash}p{0.08\textwidth} >{\raggedright\arraybackslash}p{0.31\textwidth} >{\raggedright\arraybackslash}p{0.35\textwidth} >{\raggedright\arraybackslash}p{0.18\textwidth}}
\caption{Requerimientos no funcionales}\label{tab:req-no-funcionales}\\
\hline
\textbf{ID} & \textbf{Requerimiento} & \textbf{Criterio de aceptación} & \textbf{Prioridad / estado} \\
\hline
\endfirsthead
\caption[]{Requerimientos no funcionales (continuación)}\\
\hline
\textbf{ID} & \textbf{Requerimiento} & \textbf{Criterio de aceptación} & \textbf{Prioridad / estado} \\
\hline
\endhead
\hline
\endfoot
RNF-01 & El sistema no debe emitir diagnóstico autónomo ni recomendar tratamiento. & Textos, contratos e interfaces presentan resultados como apoyo revisable e incluyen advertencias de uso no diagnóstico. & Prioridad: Alta Estado: Implementado de forma transversal \\
RNF-02 & El flujo debe ser simple y no agregar pasos innecesarios. & El recorrido principal puede completarse mediante una secuencia comprensible de autenticación, estudio, carga, análisis, visualización y revisión. & Prioridad: Alta Estado: Parcial; requiere validación de usabilidad \\
RNF-03 & El desarrollo debe ser reproducible sobre datos públicos o documentados. & Se registran dataset, versión, configuración, checkpoints, manifests, hashes, scripts, commits y evidencia de ejecución. & Prioridad: Alta Estado: Implementado con revisión continua \\
RNF-04 & El sistema debe preservar trazabilidad y auditabilidad. & Una salida puede vincularse con estudio, plano, corrida, modelo, artifact, input, corte, medición, asset y revisión. & Prioridad: Alta Estado: Implementado \\
RNF-05 & El sistema debe degradarse de forma explícita ante indisponibilidad del AI Module. & La interfaz y la API distinguen servicio disponible, degradado y error; no presentan fallback contractual como inferencia real. & Prioridad: Alta Estado: Parcial \\
RNF-06 & La comunicación y el acceso deben aplicar controles de seguridad. & Las APIs protegidas exigen autenticación, autorización por rol y manejo seguro de secretos y variables de entorno. & Prioridad: Alta Estado: Consolidado para el checkpoint P10.9; hardening productivo posterior. \\
RNF-07 & Los datos personales del paciente deben protegerse conforme a la Ley N.º 25.326. & La identidad del paciente se almacena de forma segura y con acceso restringido; los estudios se vinculan mediante una referencia seudonimizada (subjectRef). & Prioridad: Alta Estado: Implementado como criterio de diseño \\
RNF-08 & La solución debe ser modular y mantenible. & Frontend, Backend y AI Module conservan responsabilidades separadas, contratos explícitos y pruebas por repositorio. & Prioridad: Alta Estado: Implementado \\
RNF-09 & Los contratos deben ser interoperables y versionados. & Los cambios incompatibles se identifican mediante versión, fixtures y estrategia legacy; no se dependen de paths internos. & Prioridad: Alta Estado: En curso \\
RNF-10 & El rendimiento debe ser compatible con un flujo interactivo y con límites controlados. & Las cargas, timeouts y tamaños máximos se configuran; el usuario recibe estado de procesamiento y no queda ante una espera silenciosa. & Prioridad: Media Estado: Parcial \\
RNF-11 & Los resultados persistidos deben ser durables y recuperables. & Una reapertura recupera contratos y assets almacenados; la pérdida temporal del AI Module no elimina resultados existentes. & Prioridad: Alta Estado: Implementado para assets actuales \\
RNF-12 & La solución debe ofrecer observabilidad suficiente para diagnosticar fallas. & Logs, traceId, health, models, estados de ejecución y errores permiten localizar el componente responsable sin exponer datos sensibles. & Prioridad: Alta Estado: Consolidado para el checkpoint P10.9; observabilidad productiva posterior. \\
\hline
\end{longtable}
}

\section{Requerimientos de validación}

{\footnotesize
\setlength{\LTleft}{0pt}
\setlength{\LTright}{0pt}
\setlength{\tabcolsep}{2pt}
\renewcommand{\arraystretch}{1.18}
\begin{longtable}{>{\raggedright\arraybackslash}p{0.08\textwidth} >{\raggedright\arraybackslash}p{0.33\textwidth} >{\raggedright\arraybackslash}p{0.36\textwidth} >{\raggedright\arraybackslash}p{0.15\textwidth}}
\caption{Requerimientos de validación}\label{tab:req-validacion}\\
\hline
\textbf{ID} & \textbf{Requerimiento} & \textbf{Criterio de aceptación} & \textbf{Estado} \\
\hline
\endfirsthead
\caption[]{Requerimientos de validación (continuación)}\\
\hline
\textbf{ID} & \textbf{Requerimiento} & \textbf{Criterio de aceptación} & \textbf{Estado} \\
\hline
\endhead
\hline
\endfoot
RV-01 & Evaluar segmentación con métricas técnicas. & Se calculan Dice, IoU y métricas por clase o equivalentes sobre conjuntos de prueba documentados. & Implementado experimentalmente; requiere consolidación final \\
RV-02 & Complementar métricas con revisión cualitativa profesional. & Profesionales revisan ejemplos, flujo, overlays, mediciones y limitaciones; sus observaciones quedan documentadas. & En curso \\
RV-03 & Validar usabilidad y recorrido end-to-end del producto actualizado. & Se ejecuta un protocolo que cubre login, carga, análisis, persistencia, reapertura, degradación y revisión, registrando incidencias y evidencia. & En curso \\
\hline
\end{longtable}
}

\section{Matriz de trazabilidad resumida}

{\scriptsize
\setlength{\LTleft}{0pt}
\setlength{\LTright}{0pt}
\setlength{\tabcolsep}{2pt}
\renewcommand{\arraystretch}{1.18}
\begin{longtable}{>{\raggedright\arraybackslash}p{0.10\textwidth} >{\raggedright\arraybackslash}p{0.15\textwidth} >{\raggedright\arraybackslash}p{0.25\textwidth} >{\raggedright\arraybackslash}p{0.18\textwidth} >{\raggedright\arraybackslash}p{0.24\textwidth}}
\caption{Matriz de trazabilidad resumida}\label{tab:req-trazabilidad}\\
\hline
\textbf{Historia} & \textbf{Necesidad} & \textbf{Requerimientos asociados} & \textbf{Componentes} & \textbf{Evidencia esperada} \\
\hline
\endfirsthead
\caption[]{Matriz de trazabilidad resumida (continuación)}\\
\hline
\textbf{Historia} & \textbf{Necesidad} & \textbf{Requerimientos asociados} & \textbf{Componentes} & \textbf{Evidencia esperada} \\
\hline
\endhead
\hline
\endfoot
HU-01 & N-05, seguridad & RF-11, RF-12; RNF-06 & Frontend + Backend & Tests de autenticación, renovación, logout y autorización \\
HU-02 & N-07 & RF-13, RF-17 & Frontend + Backend + PostgreSQL & Captura de worklist y prueba de consulta persistida \\
HU-03 & N-10, privacidad & RF-14; RNF-07 & Frontend + Backend + PostgreSQL & Caso creado con subjectRef y sin datos identificatorios obligatorios \\
HU-04 & N-03 & RF-01, RF-07, RF-15 & Frontend + Backend + AI Module & Carga de estudio DICOM completo, identificación/registro de series compatibles y corrida multiplanar registrada. \\
HU-05 & N-03, N-06 & RF-15, RF-16; RNF-03, RNF-04 & Backend + AI Module & Contrato con runId, traceId, modelo, hash y modo efectivo \\
HU-06 & N-01 & RF-02, RF-20 & Frontend + Backend + AI Module & Overlay autorizado sobre la imagen correspondiente \\
HU-07 & N-12 & RF-09, RF-18, RF-21; RNF-09, RNF-11 & Tres repositorios & Catálogo de cortes y visor navegable (full-series); integración final P10.9 del Frontend pendiente. \\
HU-08 & N-04 & RF-04, RF-05, RF-16 & Frontend + Backend + AI Module & Medición vinculada con estructura, plano, corte y asset \\
HU-09 & N-02, N-05, N-08 & RF-03, RF-06; RNF-01, RNF-04 & Frontend + Backend & Revisión persistida con usuario, estado y observación \\
HU-10 & N-07 & RF-17, RF-18; RNF-11 & Frontend + Backend + storage & Reapertura del estudio con mediciones, correcciones y estado de revisión, validada sobre assets persistidos. \\
HU-11 & N-11 & RF-10, RF-18, RF-19; RNF-05, RNF-12 & Tres repositorios & Estado degradado comprensible y logs correlacionables \\
HU-12 & Operación técnica & RF-19; RNF-12 & Frontend + Backend + AI Module & Pantalla o endpoint protegido de salud y modelos \\
HU-13 & Seguridad & RF-12; RNF-06 & Frontend + Backend & Activación de cuenta y control de roles \\
HU-14 & N-10 & RF-08, RF-14; RNF-07, RNF-09 & Backend + datos + futura IA & Especificación longitudinal y deltas descriptivos revisables \\
\hline
\end{longtable}
}

\section{Control de alcance y requerimientos diferidos}

Las siguientes capacidades no se consideran requerimientos terminados del MVP inmediato:

\begin{itemize}
  \item diagnóstico diferencial automático;
  \item probabilidad clínica de lesión o degradación;
  \item inferencia automática sobre series o modalidades no soportadas (los cortes correspondientes se muestran como imagen base);
  \item crosshair sagital–axial sin geometría compatible;
  \item edición avanzada de máscaras comparable con un visor clínico certificado;
  \item integración productiva con PACS, RIS o historia clínica;
  \item comparación longitudinal con scoring clínico no calibrado;
  \item uso asistencial real o certificación regulatoria.
\end{itemize}

Una función propuesta que no pueda relacionarse con necesidad, historia, criterio y evidencia debe considerarse candidata a postergación. Esta regla evita que capacidades visualmente atractivas desplacen el cierre de seguridad, persistencia, navegación volumétrica y validación.

\section{Conclusión del capítulo}

La especificación conecta cada requerimiento con una necesidad del relevamiento, un componente responsable, un criterio de aceptación y su evidencia. Los requerimientos declarados como parciales o futuros se identifican como tales, de modo que la matriz de trazabilidad describe el estado real del producto y no su estado deseado.
